\section{Stakeholder und wiederkehrende Informationsbedarfe}
\label{sec:stakeholder}

Ein analytisches Informationssystem muss unterschiedliche operative,
strategische und analytische Informationsbedarfe unterstützen.
Tabelle~\ref{tab:stakeholder} fasst die zentralen internen Stakeholder und
deren wiederkehrende Anforderungen zusammen.

\begin{table}[H]
  \centering
  \caption{Stakeholder und analytische Informationsbedarfe}
  \label{tab:stakeholder}
  \begin{tabularx}{\textwidth}{@{}p{4.5cm}X@{}}
    \toprule
    Stakeholder & Analytischer Schwerpunkt und Informationsbedarf \\
    \midrule
    C-Level und Vorstand
      & Aggregierte konzernweite Kennzahlen, Geschäftsvolumen nach Segment und
        Region sowie Solvenzkennzahlen \\
    Schadenmanagement
      & Operative Schadenbearbeitung und Hinweise auf potenziell
        betrügerische Schadenfälle \\
    Vertrieb und Marketing
      & Steuerung des Neugeschäfts, Kanal-Performance,
        Zielgruppensegmentierung und Retentionsinformationen \\
    Underwriting und Aktuariat
      & Risikoeinschätzung, Tarifierungsgrundlagen und Ermittlung
        risikoadäquater Prämien \\
    Data-Science-Team
      & Rollenbasierter Zugriff auf freigegebene Roh- und Historiendaten für
        Entwicklung, Training und Überwachung von ML-Modellen \\
    \bottomrule
  \end{tabularx}
\end{table}
